\subsection{Sanity Check}
In this work, we perform 200-way zero-shot tasks.
Here, we provide an experimental chance by shuffling pairs to compare with the ideal chance of 0.5\% in Table~\ref{tab:4}.
The experimental chance is close to the ideal chance, and is significantly lower than the performance of the NICE model (\textit{p}\ \textless\ 0.001). 
% In this study, we conducted 200-way zero-shot tasks and established an experimental chance level, as presented in Table~\ref{tab:1}, for comparison with the ideal chance level of 0.5\%. The experimental chance level closely aligns with the ideal chance level and is notably lower than the performance achieved by the NICE model (\textit{p}\ \textless\ 0.001).
\label{sec:appendA0}
\begin{table}[ht]
  \caption{Chance level of the experiments}
  \label{tab:4}
  \small
  % \scriptsize
  \centering
\begin{tabular}{lcc}
   \toprule
   Condition & Top-1 (\%) & Top-5 (\%) \\
   \midrule
   ideal chance & 0.5 & 2.5 \\
   experimental chance & 0.3 & 2.3 \\
   NICE & 13.8 & 39.5 \\
   % in concepts & 11.5 & 34.9 \\
   \bottomrule
\end{tabular}
\end{table}

\subsection{Better encoders}
\label{sec:appendA1}

The adaptability of the NICE framework extends to any other well-designed EEG or image encoders, as shown in Table~\ref{tab:3}.
We further present results utilizing EVA-CLIP \cite{sun2023evaclip} pre-trained with LAION-2B as the image encoder in Table~\ref{tab:5}. The enhancements in both overall top-1 (\textit{p} < 0.001) and top-5 (\textit{p} < 0.01) accuracy are surprising.
The results further show the potential of this task and hint at the consistency of large pre-trained models and brain responses.
Beyond that, we recommend not solely focusing on overall performance, but the specific insight gained.
\begin{table}[ht]
  \caption{Overall accuracy (\%) with EVA-CLIP image encoder: top-1 and top-5}
  \label{tab:5}
  \centering
  \Huge
\resizebox{\linewidth}{!}{
  \begin{tabular}{lcccccccccccccccccccccc}
    \toprule
    & \multicolumn{2}{c}{Subject 1} & \multicolumn{2}{c}{Subject 2} & \multicolumn{2}{c}{Subject 3} & \multicolumn{2}{c}{Subject 4} & \multicolumn{2}{c}{Subject 5} & \multicolumn{2}{c}{Subject 6} & \multicolumn{2}{c}{Subject 7} & \multicolumn{2}{c}{Subject 8} & \multicolumn{2}{c}{Subject 9} & \multicolumn{2}{c}{Subject 10} & \multicolumn{2}{c}{Ave} \\
    \cmidrule(r){2-3} \cmidrule(r){4-5} \cmidrule(r){6-7} \cmidrule(r){8-9} \cmidrule(r){10-11} \cmidrule(r){12-13} \cmidrule(r){14-15} \cmidrule(r){16-17} \cmidrule(r){18-19} \cmidrule(r){20-21} \cmidrule(r){22-23}
    % \cmidrule{2-23} 
    Method & top-1 & top-5 & top-1 & top-5 & top-1 & top-5 & top-1 & top-5 & top-1 & top-5 & top-1 & top-5 & top-1 & top-5 & top-1 & top-5 & top-1 & top-5 & top-1 & top-5 & top-1 & top-5 \\
    \midrule
    \textbf{NICE} & 17.8 & 44.2 & 14.6 & 39.4 & 19.1 & 48.9 & 20.8 & 52.2 & 11.3 & 36.9 & 19.1 & 48.7 & 21.0 & 49.5 & 26.9 & 59.5 & 16.0 & 45.2 & 20.3 & 51.7 & 18.7 & 47.6 \\  
    \textbf{NICE - SA} & 17.1 & 45.4 & 15.3 & 44.2 & 17.4 & 49.0 & 22.1 & 52.0 & 15.6 & 42.4 & 21.1 & 50.5 & 23.4 & 52.5 & 26.3 & 58.8 & 18.3 & 49.1 & 22.1 & 53.5 & 19.9 & 49.7 \\
    \textbf{NICE - GA} & 19.8 & 46.7 & 19.6 & 44.7 & 19.8 & 51.5 & 25.7 & 54.5 & 13.2 & 39.2 & 23.4 & 54.1 & 22.7 & 52.9 & 26.8 & 58.9 & 17.8 & 47.1 & 23.6 & 53.9 & \multicolumn{1}{>{\columncolor{green!10}}c}{21.2} & \multicolumn{1}{>{\columncolor{green!10}}c}{50.4} \\
    \bottomrule
  \end{tabular}}
\end{table}

\subsection{Comparison of Pre-training}
\label{sec:appendA2}
We employed pre-trained models as the image encoder for two reasons: 
1) Pre-trained models like CLIP and ViT possess extensive feature extraction capabilities gained from large-scale training. 
We sought to leverage these models to enhance the EEG model, trained on a smaller dataset. 
2) Training the image encoder demands substantial computational resources. 
In our current pre-trained setup, one GPU with 8GB memory yields a well-trained EEG encoder for one subject in 5 minutes.

We conducted additional experiments without pre-trained image encoders in Table~\ref{tab:6}, training the randomly initialed image encoder and the EEG encoder simultaneously.
ViT-B/16 and ResNet-50 were used as examples. 
CLIP also relies on ViT-L, but with more layers than ViT-B (24 vs. 12), impeding completion. 
In the scenario with non-pre-trained, we observed that although training loss reduced rapidly, validation loss demonstrated a weaker decline, which is a sign of limited generalization. 
During testing, pre-trained ViT outperformed (\textit{p} < 0.01) non-pre-trained models. 
Pre-trained ResNet decreased by 1.9\% (\textit{p} < 0.05) instead, hinting at avenues of refining the image encoder to enhance overall performance.
But still, training the image encoder imposes a significant computational load. 
We should mention that the decrease in computational cost actually comes from the freezing of the pre-trained image encoder and getting image features beforehand. 

% Fine-tuning the pre-trained image model while training the EEG encoder may be another way to improve the overall performance.

\begin{table}[h]
  \caption{Comparison of pre-trained and non-pre-trained image encoders}
  \label{tab:6}
  \small
  % \footnotesize
  % \scriptsize
  \centering
  \begin{threeparttable}
\begin{tabular}{lccccc}
   \toprule
   Image Enc & Top-1 (\%) & Top-5 (\%) & Params. & Memory & Time per subject \\
   \midrule
   ViT & 7.9 & 18.5 & 88.0 M & $\sim$154 G & $\sim$ 130 min \\
   ViT* (frozen) & 12.1 & 29.4 & 1.8 M & $\sim$8 G & $\sim$5 min \\
   ResNet & 8.8 & 17.3 & 25.3 M & $\sim$112 G & $\sim$60 min \\
   ResNet* (frozen) & 6.9 & 11.8 & 1.8 M &  $\sim$8 G & $\sim$5 min \\
   \bottomrule
\end{tabular}
      \begin{tablenotes}
        \item[*]pre-trained ResNet-50 and ViT-B/16. The image features were obtained before training.  
    \end{tablenotes}
\end{threeparttable}
\end{table}

\subsection{Spatial Module Interpretation}
\label{sec:appendA3}
\begin{figure}[h]
  \centering
\includegraphics[width=1\linewidth]{Fig/Fig5.png}
  \caption{(A) Grad-CAM for self-attention (SA) and graph attention (GA) modules of individual subjects and the average across subjects. 
  (B) Attention weights of SA and GA with individual subjects and the average across subjects. 
  The visualizations highlight the regions of interest, focusing on the temporal and occipital brain areas, which are known to be associated with visual processing.}
  \label{fig:5}
\end{figure}
% \subsubsection{Features Captured by Spatial Modules}
We provide Grad-CAMs and attention weights of all subjects in Fig.~\ref{fig:5}. 
These visualizations show the alleviation of prefrontal region influence—often tied to cognitive behavior like workload and memory—by both SA and GA. 
Intriguingly, GA focuses on the occipital region, integral to visual processing, across multiple subjects. 
SA appears adept at capturing temporal region activities associated with high-level visual processing. 
% In this way, we use the model interpretation to reversely illustrate that EEG responses can perceive brain activity on occipital and temporal areas, associated with object recognition.
Our model interpretation demonstrates that EEG responses can capture and reveal brain activity in the occipital and temporal areas, which are strongly associated with object recognition processes. 
This provides implicit evidence for the effectiveness of EEG-based object recognition, and further validates the plausibility of our approach.

\subsection{Parameter Sensitivity}
\label{sec:appendA4}
The temperature parameter (\(\tau\)) in contrastive learning plays a crucial role in shaping the similarity distribution. 
In our implementation, we optimized this temperature directly as a log-parameterized multiplicative scalar during training, similar to the CLIP model. 
To explore the impact of temperature settings, we conducted a series of comparisons, and the results are presented in Table~\ref{tab:7}. These comparisons demonstrate the significant influence of temperature adjustments on the results.
Besides, our pre-experiments help determine some conventional parameters such as batch size, decay, etc.

\begin{table}[h]
  \caption{Comparison of the temperature parameter in contrastive learning}
  \label{tab:7}
  \small
  % \footnotesize
  \centering
\begin{tabular}{lccccccccc}
   \toprule
   Temperature & learnable & 0.001 & 0.005 & 0.01 & 0.03 & 0.05 & 0.1 & 0.5 \\
   \midrule
   Top-1 (\%) & 13.8 & 12.4 & 13.45 & 15.8 & 15.2 & 14.7 & 12.1 & 10.5 \\
   signif. (p) & - & >.05 & >.05 & <.01 & <.01 & <.05 & <.01	& <.001 \\
   \bottomrule
\end{tabular}
\end{table}

\subsection{MEG Experiments}
\label{sec:appendA5}

We provided analysis on an additional MEG dataset \cite{hebart2023thingsdata} to corroborate the results on EEG data.
The MEG dataset of four participants has 271 channels and more stable responses with a long stimulus duration of 500 ms followed by a blank screen of 1000\(\pm\)200 ms.
There are 1854 concepts\(\times\)12 images\(\times\)1 repetitions in the training stage and 200 concepts\(\times\)1 image\(\times\)12 repetitions in the test stage.
We discarded 200 test concepts from the training set to construct the zero-shot task.
The MEG data were epoched into trials from 0 to 1000 ms after the stimuli onset.
We used a band-pass filter of [0.1, 100] Hz and baseline correction after down-sampling to 200 Hz for preprocessing.
We averaged all MEG repetitions of one image to ensure the signal-to-noise ratio.
Note that the statistical analysis was not performed on the MEG dataset due to the few participants.
\begin{table}[htpb]
\caption{Overall accuracy (\%) of 200-way zero-shot classification and retrieval on MEG data}
\label{tab:8}
\small
% \footnotesize
\centering
\resizebox{\linewidth}{!}{
\begin{tabular}{clcccccccccc}
    \toprule
    & & \multicolumn{2}{c}{Sub 1} & \multicolumn{2}{c}{Sub 2} & \multicolumn{2}{c}{Sub 3} & \multicolumn{2}{c}{Sub 4} & \multicolumn{2}{c}{Ave} \\
    \cmidrule(r){3-4} \cmidrule(r){5-6} \cmidrule(r){7-8} \cmidrule(r){9-10} \cmidrule(r){11-12} 
    % \cmidrule{2-23} 
    Task & Method & Top-1 & Top-5 & Top-1 & Top-5 & Top-1 & Top-5 & Top-1 & Top-5 & Top-1 & Top-5 \\
    \midrule
    \multirow{3}{*}{classification} 
    & NICE & 6.9 & 20.5 & 15.3 & 37.1 & 12.3 & 35.0 & 5.8 & 21.1 & \textbf{10.1} & \textbf{28.4} \\
    & w/ SA & 7.3 & 22.6 & 13.1 & 37.1 & 10.2 & 30.2 & 8.3 & 24.5 & 9.7 & 28.6 \\ 
    & w/ GA & 6.4 & 23.4 & 16.2 & 40.8 & 14.0 & 38.9 & 6.4 & 25.2 & 10.8 & 32.1 \\
    \midrule
    \multirow{3}{*}{retrieval} 
    & NICE & 9.6 & 27.8 & 18.5 & 47.8 & 14.2 & 41.6 & 9.0 & 26.6 & \textbf{12.8} & \textbf{36.0} \\
    & w/ SA & 9.8 & 28.1 & 18.6 & 46.4 & 10.5 & 38.4 & 11.7 & 27.2 & 12.7 & 35.0 \\
    & w/ GA & 8.7 & 30.5 & 21.8 & 56.6 & 16.5 & 49.7 & 10.3 & 32.3 & 14.3 & 42.3 \\
    \bottomrule 
\end{tabular}}
\end{table}

\begin{figure}[htbp]
  \centering
  \includegraphics[width=1\linewidth]{Fig/Fig6.png}
  \caption{Temporal, spatial, spectral analysis of MEG data. 
  (A) Topographies of each 100 ms by averaging training trials. 
  (B) Averaged accuracy with different time lengths.
  (C) Averaged accuracy of different brain regions (not ablation here for clarity).
  (D) Time-frequency maps of the occipital, temporal, and parietal data from one subject. 
  (E) Averaged accuracy of different rhythms.}
  \label{fig:6}
\end{figure}	

% \subsubsection{Classification and Retrieval}
% \subsubsection{Temporal, Spatial, and Spectral Dynamics}
Classification and retrieval tasks with the MEG dataset have been conducted with the NICE framework in Table~\ref{tab:8}. 
A similar 200-way zero-shot classification showed above-chance performance on MEG data with an average top-1 accuracy of 10.1\% and top-5 accuracy of 28.4\%.
Besides, we provided the results on a retrieval task, which meant directly using the stimulus images for matching templates instead of other images belonging to the concepts.
The spatial module GA could still help with the results, but SA varied across individuals.
It is worth noting that a more suitable encoder for extracting MEG features should help with the performance \cite{benchetrit2023brain}.
Moreover, we roughly conducted the temporal, spatial, and spectral analysis of MEG data in Fig.~\ref{fig:6}.
The results were similar to Fig.~\ref{fig:2}, while the temporal responses changed because of the longer stimuli duration.

\subsection{EEG Encoder Comparison}
\label{sec:appendA6}
The residual connection \cite{he2016deepa} was used with SA and GA modules to stabilize the training procedure. 
We compared the effect of the residual connection in Table~\ref{tab:9}. 
It can be seen that it significantly helped GA in top-1 accuracy (p<0.01) but no significance for SA (p > 0.05). 

The EEG encoder TSConv in our framework has a concise architecture consisting of basic temporal and spatial convolutional layers.
We compared the details in Table~\ref{tab:10} because the design was similar to ShallowNet \cite{schirrmeister2017deepa}.
The comparison encompassed several issues, including the kernel size, pooling order, and activation function.
All three cases caused a significant effect on the top-1 accuracy compared to TSConv (p < 0.01). 
Besides, we provided the number of parameters and the output dimension within different EEG encoders in Table~\ref{tab:11}.
We believe that other well-designed EEG feature extractors can help us improve the performance in this framework.
\begin{figure}[b]
  \centering
  \includegraphics[width=1\linewidth]{Fig/Fig7.png}
  \caption{t-SNE visualization of four categories: animal, food, vehicle, and tool. 
  (A) Visual feature distribution of the training and test sets.
  (B) EEG feature distribution of the training and test sets.}
  \label{fig:7}
\end{figure}	

\begin{table}[htpb]
\caption{Residual connection comparison for spatial modules}
\label{tab:9}
\small
% \footnotesize
\centering
\begin{tabular}{clcc}
    \toprule
    Module & Res & Top-1 (std) & Top-5 (std) \\
    \midrule
    \multirow{2}{*}{SA} 
    & w/ res & 14.7 (2.3) & 41.7 (4.6) \\
    & w/o res & 13.8 (2.8) & 39.6 (5.7) \\ 
    \midrule
    \multirow{2}{*}{GA} 
    & w/ res & 15.6 (2.9) & 42.8 (5.5) \\
    & w/o res & 10.1 (3.2) & 29.9 (5.3) \\
    \bottomrule 
\end{tabular}
\end{table}

\begin{table}[htpb]
\caption{Comparison of the architecture of TSConv}
\label{tab:10}
\small
\centering
\begin{tabular}{lcc}
    \toprule
    Condition & Top-1 (std) & Top-5 (std) \\
    \midrule
    original ShallowNet & 7.9 (4.3) & 24.2 (9.3) \\
    kernel size as ShallowNet & 9.9 (2.4) & 33.5 (5.4) \\ 
    pooling order & 10.9 (3.0) & 34.3 (7.2) \\
    ELU->ReLU & 11.0 (2.5) & 33.3 (6.6) \\
    TSConv & 13.8 (3.3) & 39.5 (6.5) \\
    \bottomrule 
\end{tabular}
\end{table}

\begin{table}[htpb]
\caption{Details of different EEG encoders}
\label{tab:11}
\small
\centering
\begin{tabular}{lcc}
    \toprule
    EEG encoder & Params & Out dim \\
    \midrule
    ShallowNet & 102.0 k & 440 \\
    DeepNet & 303.3 k & 1400 \\ 
    Conformer & 161.2 k & 1440 \\
    EEGNet & 12.8 k & 1248 \\
    TSConv & 102.0 k & 1440 \\
    \bottomrule 
\end{tabular}
\end{table}

\subsection{Features Distribution}
\label{sec:appendA7}
We plotted t-SNE \cite{maaten2008visualizinga} in Fig.~\ref{fig:7} to show the distribution variance between training and test sets with visual and EEG features separately, to be a supplement of the RSA in Fig.~\ref{fig:3}.
Four categories, animal, food, vehicle, and tool, were selected for the visualization, which revealed why we could achieve zero-shot performance after training with rich data. 
